% TOP_PROBLEM: {"id": 30006700, "title": "Equivariant Tamagawa number conjecture: Burns-Flach Conjecture4 package", "category_id": 1, "category": {"id": 1, "name": "number_theory", "display_name": "Number Theory", "description": "Properties of integers, prime numbers, Diophantine equations.", "slug": "number-theory", "order_index": 1, "created_at": "2026-07-31T15:26:25.670Z"}, "status": "open"}
% FIRST_POSTED: 2026-09-27
% !TeX program = lualatex
% Compile twice: lualatex 30006700.tex
% Original entry retained; research subsection and [E] references added.
% No external bibliography, image, data, or style file is required.
\documentclass[11pt,a4paper]{article}
\usepackage[margin=24mm,headheight=14pt]{geometry}
\usepackage{fontspec}
\setmainfont{Latin Modern Roman}
\setsansfont{Latin Modern Sans}
\setmonofont{Latin Modern Mono}
\usepackage{amsmath,amssymb,mathtools,amsthm}
\usepackage{unicode-math}
\setmathfont{Latin Modern Math}
\usepackage{microtype}
\usepackage{enumitem,xurl,needspace}
\usepackage[dvipsnames]{xcolor}
\usepackage{hyperref}
\hypersetup{unicode=true,colorlinks=true,linkcolor=MidnightBlue,
 urlcolor=MidnightBlue,bookmarksnumbered=true,
 pdftitle={Research notebook - Problem 30006700},
 pdfauthor={Research attempt; not independently refereed}}
\usepackage{bookmark}
\usepackage{titlesec}
\titleformat{\section}{\Large\bfseries\raggedright}{\thesection}{0.7em}{}
\usepackage{fancyhdr}
\pagestyle{fancy}\fancyhf{}
\fancyhead[L]{\small\sffamily Research notebook}
\fancyhead[R]{\small\sffamily Problem 30006700}
\fancyfoot[L]{\scriptsize 27 September 2026}
\fancyfoot[C]{\thepage}
\fancyfoot[R]{\scriptsize Unrefereed research attempt}
\setlength{\parindent}{0pt}
\setlength{\parskip}{5pt plus 1pt}
\setlength{\emergencystretch}{3em}
\clubpenalty=10000
\widowpenalty=10000
\displaywidowpenalty=10000
\setcounter{secnumdepth}{1}
\setlist[itemize]{leftmargin=3em,itemsep=5pt,topsep=4pt,parsep=0pt}
\allowdisplaybreaks
\newcommand{\N}{\mathbb{N}}
\newcommand{\Z}{\mathbb{Z}}
\newcommand{\Q}{\mathbb{Q}}
\newcommand{\R}{\mathbb{R}}
\newcommand{\C}{\mathbb{C}}
\newcommand{\F}{\mathbb{F}}
\newcommand{\A}{\mathbb{A}}
\newcommand{\PP}{\mathbb P}
\newcommand{\E}{\mathbb{E}}
\newcommand{\eps}{\varepsilon}
\newcommand{\dd}{\,\mathrm d}
\newcommand{\sref}[2]{\hyperlink{source:#1:#2}{[S#2]}}
\newcommand{\eref}[2]{\hyperlink{extra:#1:#2}{[E#2]}}
\newif\ifshowcatalogue
\showcataloguetrue
\newcommand{\cataloguescope}[1]{\ifshowcatalogue
 \paragraph{Exact scope in the supplied catalogue.}
 \begin{quote}\small #1\end{quote}\fi}
\newtheorem{notebooklemma}{Lemma}[section]
\newtheorem{notebookproposition}[notebooklemma]{Proposition}
\numberwithin{equation}{section}
\begin{document}
\setcounter{section}{0}
% ================================================================
% problemId: 30006700
% Canonical problem ID: 30006700
% exactTarget SHA-256: cb6199835108f5143aba0756565afe9f0debb0b815f2153696c92827ed695da8
\clearpage
\section*{\texorpdfstring{Equivariant Tamagawa number conjecture: Burns-Flach Conjecture4 package}{Equivariant Tamagawa number conjecture: Burns-Flach Conjecture4 package}}\label{30006700}
\subsection{Definitions and mathematical statement}
Let $A$ be a finite-dimensional semisimple $\Q$-algebra and $\mathcal A\subset A$ an order with the compatible projective structure on $M$ specified by Burns--Flach. Relative $K_0(\mathcal A,\R)$ records projective modules together with isomorphisms after extension to $A\otimes\R$. The conjecture is the \emph{entire} conditional package of Conjecture 4(i)--(iv): meromorphic continuation at zero, the reduced-rank order formula, rationality of the Tamagawa class, and its vanishing. In particular,
\[
 r=\operatorname{rr}_A H_f^1(K,M^\vee(1))^*
   -\operatorname{rr}_A H_f^0(K,M^\vee(1))^*,
\]
\[
 T_\Omega=\widehat\delta^1_{\mathcal A,\R}\bigl(L^*(AM,0)\bigr)
             +R_\Omega(M,\mathcal A)=0.
\]
Here $L^*$ is the componentwise leading Laurent coefficient, $R_\Omega$ is the fundamental virtual-object class with its comparison trivialization, and $\widehat\delta$ is the canonical extended reduced-norm boundary. Their exact definitions and the external comparison, finiteness, coefficient-compatibility, and coherence hypotheses are incorporated from the source. Noncommutative and nonmaximal orders remain in scope. The full recorded target below prevents this last identity from being mistaken for the whole package.
\par\smallskip\textit{Formulation and concept references:} \sref{30006700}{1}.
\cataloguescope{For every motive M over a number field K in Burns-Flach Section3.1's motivic-structure formalism, with an action of a finite-dimensional semisimple Q-algebra A, and every order mathcal A in A admitting the compatible projective mathcal A-structure of Section3.3, prove Burns-Flach Conjecture4(i)-(iv) in its stated conditional setting. The external defining hypotheses are the finite motivic cohomology and real comparison exact sequence of Conjecture1, the motivic-to-p-adic comparison of Conjecture2 for M and Mdual(1), coefficient compatibility giving well-defined equivariant local factors as in4.1 (Conjecture3 or the weaker Euler-factor compatibility of Remark7), and the stated Coherence hypothesis. A projective structure consists of full projective lattices in the archimedean realizations whose l-adic comparison images agree and are Galois-stable, not merely an arbitrary maximal-order choice. The INCLUDED assertions are: (i) the finite equivariant L-function L(AM,s) has meromorphic continuation to0; (ii) its componentwise integer order r equals rr\_A(H\_f\textasciicircum{}1(K,Mdual(1))\textasciicircum{}*)-rr\_A(H\_f\textasciicircum{}0(K,Mdual(1))\textasciicircum{}*), with reduced rank as in2.6; (iii) for the nonzero componentwise leading Laurent coefficient L*(AM,0)=lim\_(s->0) s\textasciicircum{}(-r)L(AM,s), the class T\_Omega=delta\_hat\textasciicircum{}1\_(mathcal A,R)(L*(AM,0))+R\_Omega(M,mathcal A) lies in Cl(mathcal A,Q); and (iv) T\_Omega=0 in Cl(mathcal A,R), hence in K\_0(mathcal A,R). The R\_Omega term is the class of the fundamental virtual object with its integral compact-support cohomology structures and real comparison trivialization in3.4; delta\_hat is Lemma9's canonical extended reduced-norm boundary, not an inverse to a surjective real reduced norm. Equivalently the final identity is delta\_hat(L*)=-R\_Omega. Preserve graded determinant conventions and all suitable nonmaximal/noncommutative orders. Coherence-free Remark8's Conjecture5 and Conjecture6 for EVERY prime are an alternative to the rational/integral clauses in the source's stated sense, not removal of(i)-(ii) or replacement by a single p-local or epsilon assertion. The external hypotheses are not presumed universally proved or restricted to currently established cases, and proving all of them is not added as a separate conclusion.}
\subsection{Short English statement}
Do equivariant motivic $L$-functions satisfy the full Burns--Flach continuation, vanishing-order, rationality, and integral leading-term law under its stated prerequisites?
\subsection{Research attempt}

\paragraph{Current frontier.}
The source-defined target is the full Burns--Flach package, not just a special-value formula after extending scalars to a maximal order. Its external hypotheses remain assumptions, whereas continuation, the order formula, rationality, and integral vanishing remain conclusions \eref{30006700}{1}. Johnston--Nickel prove an odd-prime equivariant Iwasawa main conjecture for the specified totally real, one-dimensional extensions with abelian Sylow $p$-subgroup, obtaining applications to Tate motives \eref{30006700}{2}. Longo--Vigni's 2026 paper establishes cases of the $p$-part for higher-weight modular motives in analytic ranks zero and one, under explicit regulator and Abel--Jacobi hypotheses and further conditions on the form and prime \eref{30006700}{3}. Neither result supplies the entire all-motive, all-order package here.

\paragraph{Attack route.}
Three possible reductions must be separated. Analytic continuation and the order-of-vanishing formula might be approached through automorphy, but arbitrary motives in the source's formalism are not all known to be automorphic. Character-by-character leading-term identities might establish a maximal-order statement, but they can forget integral congruences between characters. Finally, a rational class could be checked at all primes; that is useful only after rationality has been supplied. We investigate the second and third routes together, by calculating exactly what can be lost at a nonmaximal order and identifying a finite collection of residual congruence obstructions.

\paragraph{Attempt: preserving the conditional package.}
Keep the source's realizations, finite motivic cohomology and real comparison sequence, motivic-to-$p$-adic comparisons for both $M$ and $M^\vee(1)$, coefficient compatibility, and Coherence. In particular the projective lattices must agree under comparison and be Galois-stable; an arbitrary lattice in one realization does not qualify. We retain
\[
 T_\Omega=\widehat\delta^1_{\mathcal A,\R}(L^*(AM,0))
                       +R_\Omega(M,\mathcal A),
\]
with the source's graded determinant conventions and its extended boundary, rather than an assumed inverse to a surjective real reduced norm \eref{30006700}{1}, Sections 3--4 and Lemma 9.

In the deductions below, clauses (i)--(iii) will be \emph{additional hypotheses of a reduction}, not asserted consequences of the external comparisons. In particular, no analytic continuation is manufactured by the algebraic calculation. The relevant source input, once rationality is available, is
\begin{equation}\label{r58:localization}
 \operatorname{Cl}(\mathcal A,\Q)
 \simeq\bigoplus_p
 \frac{K_1(A_p)}{\operatorname{im}K_1(\mathcal A_p)},
 \qquad A_p=A\otimes_{\Q}\Q_p,
 \quad\mathcal A_p=\mathcal A\otimes_{\Z}\Z_p,
\end{equation}
from Burns--Flach, equation (15). Its injection into the real relative group allows vanishing to be tested on these components \eref{30006700}{1}, Section 2.9.

\paragraph{Attempt: a nonmaximal-order test that defeats componentwise detection.}
Fix a prime $p$ and an integer $e\geq1$. Put
\[
 B=\Q_p\times\Q_p,\qquad
 \Gamma=\Z_p\times\Z_p,\qquad
 \Lambda_e=\{(a,b)\in\Z_p^2:a\equiv b\pmod{p^e}\}.
\]
The ring $\Lambda_e$ is the $p$-adic completion of the genuine global order
\[
 \mathcal L_e=\{(a,b)\in\Z^2:a\equiv b\pmod{p^e}\}
                           \subset\Q\times\Q.
\]
Thus the example is not an artifact of allowing rings outside the order setting.

The units of $\Lambda_e$ are exactly the pairs of $p$-adic units that are congruent modulo $p^e$. Indeed the inverse pair remains congruent because
$a^{-1}-b^{-1}=(b-a)/(ab)$. Every nonunit reduces to zero under the common residue map to $\F_p$. Hence $\Lambda_e$ is a commutative local ring, all its finitely generated projective modules are free, and
\[
 K_0(\Lambda_e)=\Z\longrightarrow K_0(B)=\Z^2,
 \qquad 1\longmapsto(1,1)
\]
is injective. For completeness, freeness follows by lifting a residue-field basis of a finitely generated projective module: Nakayama's lemma gives a surjection from a finite free module, projectivity splits it, and the kernel has zero reduction and hence is zero. Over a commutative local ring, Gaussian elimination by unit pivots identifies $K_1$ with the unit group via determinant. The relative $K$-theory exact sequence consequently gives
\begin{equation}\label{r58:toyrelative}
 K_0(\Lambda_e,\Q_p)\simeq
       \frac{(\Q_p^\times)^2}{\Lambda_e^\times}.
\end{equation}
For $\Gamma$, the analogous relative group is $\Z^2$, detected by the two valuations.

\begin{notebookproposition}[The lost unit congruence]\label{r58:kernel}
Extension of scalars has kernel
\[
 \ker\!\left(K_0(\Lambda_e,\Q_p)\longrightarrow
                  K_0(\Gamma,\Q_p)\right)
       \simeq(\Z/p^e\Z)^\times.
\]
\end{notebookproposition}
\begin{proof}
By \eqref{r58:toyrelative}, the kernel is $\Gamma^\times/\Lambda_e^\times$. The homomorphism
\[
 \Gamma^\times\longrightarrow(\Z/p^e\Z)^\times,
 \qquad(u_1,u_2)\longmapsto u_1u_2^{-1}\bmod p^e
\]
is surjective and its kernel is precisely $\Lambda_e^\times$. The asserted isomorphism follows.
\end{proof}

For $p=5$, $e=1$, multiplication by $u=(1,2)$ on $B$ gives the relative triple $(\Lambda_1,u,\Lambda_1)$. Its image in $K_0(\Gamma,\Q_5)$ is zero because both components are units. Its class over $\Lambda_1$ is nonzero: the residue ratio is $3\in\F_5^\times$, of exact order four. In particular,
\begin{equation}\label{r58:falseimplication}
 \begin{gathered}
 \text{componentwise and maximal-order vanishing}\\
 \not\Longrightarrow\ \text{vanishing over the original order}.
 \end{gathered}
\end{equation}
The relative exact sequence used here is the same one underlying the source's formalism \eref{30006700}{1}, equation (10); all the concrete group identifications have been derived above.

The conductor of $\Lambda_e$ in $\Gamma$ is $p^e\Gamma$. To verify this, multiply a candidate conductor element $(a,b)$ by $(1,0)$ and $(0,1)$; both coordinates must be divisible by $p^e$, and that condition plainly suffices. Thus the order is the pullback of the diagonal subring of $(\Z/p^e\Z)^2$. The residual ratio is exactly the unit-gluing information at this conductor, not an unexplained torsion term.

\paragraph{Attempt: a finite-prime reduction for the original class.}
Return to a general semisimple $\Q$-algebra $A$ and its order $\mathcal A$. Choose a maximal overorder $\mathcal G\supset\mathcal A$. The finite abelian group $\mathcal G/\mathcal A$ is supported on a finite set of rational primes
\[
 S=\{p:\mathcal A_p\ne\mathcal G_p\}.
\]
For $p\in S$ define an abelian group
\begin{equation}\label{r58:residual}
 D_p=\frac{\operatorname{im}(K_1(\mathcal G_p)\to K_1(A_p))}
           {\operatorname{im}(K_1(\mathcal A_p)\to K_1(A_p))}.
\end{equation}
The denominator is a subgroup of the numerator because the inclusion of orders factors through $\mathcal G_p$. We deliberately keep the full $K_1$ images: replacing them by an unsupported reduced-norm identification in the noncommutative case would be a gap.

\begin{notebookproposition}[Conductor-sensitive reduction]\label{r58:reduction}
Assume the external hypotheses of the original package and clauses (i), (ii), and (iii). Suppose the image of $T_\Omega$ over $\mathcal G$ vanishes. Then $T_\Omega$ determines a class $\eta_p\in D_p$ for each $p\in S$, and
\[
 T_\Omega=0\quad\Longleftrightarrow\quad
                \eta_p=0\text{ for every }p\in S.
\]
Consequently, clauses (i)--(iii), maximal-order vanishing, and these residual vanishing statements suffice for the \emph{full} Conjecture 4(i)--(iv).
\end{notebookproposition}
\begin{proof}
Rationality places $T_\Omega$ in \eqref{r58:localization}. Represent its $p$-component by an element $\gamma_p\in K_1(A_p)$. Its image in the maximal-order quotient is zero, so $\gamma_p$ belongs to the image of $K_1(\mathcal G_p)$. The resulting coset modulo the image of $K_1(\mathcal A_p)$ is independent of the representative and is precisely $\eta_p$. At $p\notin S$ the two orders agree, so the component is already zero. Formula~\eqref{r58:localization} now proves both directions. Clauses (i)--(iii) were assumed, not deduced from this vanishing argument.
\end{proof}

Scalar extension of a compatible projective structure is taken throughout when comparing the motivic class with its maximal-order image; it is not merely a change in the coefficient field of the $L$-function. The source treats the relevant functoriality in Section 4.4. Its Proposition 4.2 provides a commutative reduction with a \emph{maximal-order} qualification, which is exactly why \eqref{r58:falseimplication} is an appropriate stress test \eref{30006700}{1}.

In the two-component order, Proposition~\ref{r58:reduction} asks for one very concrete condition: after the componentwise valuations have been accounted for, the remaining two comparison units must be congruent modulo $p^e$. Checking the valuations of the two leading terms separately never certifies this.

\paragraph{Attempt: can congruences of families force the missing gluing?}
A potentially useful but limited observation is the following. If
$F_1(t),F_2(t)\in\Z_p[[t]]$ are congruent coefficientwise modulo $p^e$, have the same order $r$ at $t=0$, and their leading coefficients are units, then those leading coefficients are congruent modulo $p^e$. This is proved simply by comparing the coefficient of $t^r$. It suggests constructing one integral equivariant zeta object \emph{before} projecting to characters and extracting leading terms.

There are two serious obstructions to using that observation as a proof. First, the included order formula is componentwise, and component orders can differ. Second, division by comparison periods or regulators need not preserve an available congruence. For a concrete algebraic test, take $p$ odd and
\[
 F_1(t)=t,\qquad F_2(t)=t+2p^e.
\]
The series are congruent modulo $p^e$. Their orders at zero are respectively one and zero. After removing the $p$-power valuation from each nonzero leading coefficient, the remaining units are $1$ and $2$, not congruent modulo $p^e$. Thus coefficientwise congruence alone does not survive arbitrary componentwise leading-term normalization. This example models a failed inference; it is not asserted to be a pair of motivic $L$-functions.

A successful version of this route would have to construct integral comparison data over the actual order, compatible both with variable vanishing orders and with the graded determinant of compact-support cohomology. Producing such a universal object is not an established consequence of the external comparison hypotheses. Assuming it would conceal much of the desired integral assertion.

\paragraph{Stress test.}
The class represented by $(1,2)$ is \emph{not} a counterexample to the equivariant Tamagawa number conjecture. No motive has been constructed whose actual $T_\Omega$ equals that class. Arbitrarily changing one comparison map, or changing a period while leaving its lattice untouched, is not allowed by the source-defined package. Consistent basis changes leave a relative class unchanged; the source's construction also controls changes of projective structure \eref{30006700}{1}, Lemma 5.

Nor does this reduction prove the analytic clauses. The statements about continuation and reduced rank, and the rationality required to enter \eqref{r58:localization}, remain distinct unsolved inputs in the universal problem. An all-prime argument cannot be replaced by the single ordinary-prime modular results in \eref{30006700}{3}; an odd-prime Tate-motive theorem cannot be extended to all primes and motives by dropping its hypotheses \eref{30006700}{2}. The relation to catalogue Section 59 must likewise not be read as an unconditional proof of finiteness of every elliptic Tate--Shafarevich group: the motivic package itself has external conditional prerequisites.

The finite verification script computes the following unit-gluing quotient sizes:
\[
 \begin{array}{c|ccccc}
 (p,e)&(2,1)&(2,2)&(3,2)&(5,1)&(5,2)\\ \hline
 |(\Z/p^e\Z)^\times|&1&2&6&4&20.
 \end{array}
\] It also verifies the order-four class of $(1,2)$ at $(5,1)$. These are exhaustive computations in finite unit groups, not calculations of a motivic Tamagawa class. They check in particular the trivial edge case $(2,1)$; the obstruction is not falsely asserted to be nonzero for every conductor.

\paragraph{Outcome.}
\textbf{Outcome: Conditional reduction.}

The work isolates a conductor-sensitive residual condition and proves that, once clauses (i)--(iii) and maximal-order vanishing are supplied, this condition is exactly the remaining integral obstruction. A fully explicit nonmaximal order shows that componentwise vanishing does not formally imply it. What remains is an arithmetic theorem forcing the actual motivic residual classes to vanish, together with proofs of the analytic and rationality clauses in the required generality. The elementary order calculation and reduction have not been independently refereed; they are not claimed as new results in relative $K$-theory, nor as a resolution of any unproved motivic case.

\subsection{Sources}
\begin{itemize}\raggedright
\item[\textnormal{[S1]}]\hypertarget{source:30006700:1}{} David Burns and Matthias Flach, Tamagawa Numbers for Motives with (Non-Commutative) Coefficients, Documenta Mathematica6(2001),501-570.
\url{https://ems.press/content/serial-article-files/25892}.
{\small\textit{Catalogue formulation source. Consulted 24 September 2026: Sections 3--4; relative K-theory and extended boundary.}}
\item[\textnormal{[S2]}]\hypertarget{source:30006700:2}{} The Tamagawa number conjecture and Kolyvagin's conjecture for motives of modular forms.
\url{https://link.springer.com/article/10.1007/s40687-026-00646-7}.
{\small\textit{Catalogue research/status reference; not a proof endorsement.}}
\item[\textnormal{[S3]}]\hypertarget{source:30006700:3}{} An unconditional proof of the abelian equivariant Iwasawa main conjecture and applications.
\url{https://athene-forschung.unibw.de/92554?show_id=155671}.
{\small\textit{Catalogue research/status reference; not a proof endorsement.}}

% Research references added for this notebook.
\item[\textnormal{[E1]}]\hypertarget{extra:30006700:1}{} David Burns and Matthias Flach, \emph{Tamagawa Numbers for Motives with (Non-Commutative) Coefficients}, Documenta Mathematica 6 (2001), 501--570: Section 2.9, equations (10)--(15); Sections 3.1--3.4, including the comparison, projective-structure and Coherence hypotheses; Lemma 9; Conjecture 4(i)--(iv); Proposition 4.2 and Section 4.4.
\url{https://ems.press/content/serial-article-files/25892}.
{\small\textit{Detailed formulation and relative K-theory used here. The original [S1] entry is retained; this [E1] gives the precise locations actually used. Consulted 27 September 2026.}}
\item[\textnormal{[E2]}]\hypertarget{extra:30006700:2}{} Henri Johnston and Andreas Nickel, \emph{An unconditional proof of the abelian equivariant Iwasawa main conjecture and applications}, arXiv:2010.03186v3 (6 December 2024), abstract and Introduction; revision records changes following a referee report.
\url{https://arxiv.org/abs/2010.03186v3}.
{\small\textit{Odd-prime, totally real, specified one-dimensional extensions; applications include Tate-motive cases, not the full all-motive package. Consulted 27 September 2026.}}
\item[\textnormal{[E3]}]\hypertarget{extra:30006700:3}{} Matteo Longo and Stefano Vigni, \emph{The Tamagawa number conjecture and Kolyvagin's conjecture for motives of modular forms}, Research in the Mathematical Sciences 13 (2026), article 64, published 30 July 2026, abstract and Section 1. DOI: 10.1007/s40687-026-00646-7.
\url{https://link.springer.com/article/10.1007/s40687-026-00646-7}.
{\small\textit{Conditional p-part statements for modular motives; regulator and Abel--Jacobi hypotheses are not silently dropped. Consulted 27 September 2026.}}
\end{itemize}

\end{document}
